\subsection{System Overview}
This section establishes the foundational context of the aBridgeAI system, outlining its core mission, the entities invested in its success, and the human actors who interact with the platform.

\subsubsection{Purpose}
The purpose of aBridgeAI is to bridge the gap between theoretical academic curricula and practical industry expectations. By operating as an AI-driven LMS and Mock Interview system, it transforms traditional, passive learning into an active, context-aware application loop. The system aims to automate the generation of assessments and provide personalized interview practices to ensure students are "Job-Ready."

\subsubsection{Stakeholders and Users of the System}
It is essential to distinguish between stakeholders (those impacted by or invested in the system's outcome) and users (those directly interacting with the software).

\paragraph{Stakeholders}
\begin{itemize}
    \item \textbf{Higher Education Institutions and Coding Bootcamps:}
    The intended institutional adopters of aBridgeAI, which can use
    the platform to reduce assessment overhead and inspect progression
    evidence within their authorised scope.

    \item \textbf{Accreditation Bodies, Industry Partners, and Hiring Companies:}
    External stakeholders that may benefit indirectly from institutional
    progression evidence and improved graduate preparation. They are not
    direct system users, and the current platform specifies no direct
    reporting, data-sharing, or candidate-verification workflow for them.
 
    \item \textbf{Academic Supervisor and Project Review Committee:}
    The faculty advisors and reviewers responsible for evaluating
    the scientific and engineering contributions of this work
    against academic criteria.
\end{itemize}

\paragraph{Users of the System}
\begin{itemize}
    \item \textbf{Student:} The primary end-user. Consumes
    multimedia lessons, participates in AI-driven quiz and spaced-repetition quiz
    review cycles whose intervals are governed by the hybrid SM-2
    scheduler, takes AI-conducted mock interviews in a unified hybrid
    room supporting typed and spoken interaction, and tracks personal
    progression via the Knowledge Retention Estimate, Progression
    Readiness, and Review Compliance Rate metrics. The student receives
    a binary pass/fail outcome together with qualitative remediation and
    a study plan; numeric rubric scores, per-outcome verdicts, and model
    reasoning remain hidden.
 
    \item \textbf{Teacher:} Designs lessons, uploads
    raw materials for AI ingestion, and is
    responsible for reviewing AI-drafted quiz and interview questions
    before publication --- including editing quiz cards, confirming the
    per-card expected response time $T_{\text{exp}}$ that anchors $Q$
    derivation, defining target outcomes for the mock interview, and
    reviewing the AI-generated interview question bank. Material
    ingestion and derived retrieval units are processed automatically
    rather than individually teacher-approved. Teachers additionally
    configure the SM-2 unlock
    parameters ($EF_{\min}^{\text{unlock}}$, coverage threshold
    $\tau$, prerequisite criteria) per lesson and monitor cohort
    retention and compliance metrics via the teacher dashboard.
 
    \item \textbf{Educational Manager:} An academic operations role
    acting within the organisation or faculty scope granted by its
    permissions (e.g.\ a programme lead). May manage pathways,
    teaching assignments, programmes, and enrolments only within that
    resolved scope; bulk enrolment and organisation-scoped reporting
    remain subject to the same boundary. Faculty Deans retain the
    distinct authority to adjudicate pathway-change requests. An
    Educational Manager has no cross-tenant visibility.

    \item \textbf{Faculty Dean:} The academic authority over a single
    faculty or department within an organisation (e.g.\ a Dean or Vice
    Dean). The Dean holds the full authority of an Educational Manager
    within their assigned organisational unit---the dean permission
    set subsumes the manager one---and additionally owns the
    governance decisions that belong to a faculty rather than to the
    organisation as a whole. Chief among these is authoring learning
    programmes within their faculty and adjudicating student requests
    to change career pathway or to drop one: the Dean acknowledges,
    approves, or rejects each request, and an approval either transfers
    the student's earned progress onto the incoming pathway or closes
    the dropped one against a snapshot of what it had earned. The
    Dean's authority is bounded by the faculty they administer.
 
    \item \textbf{IT Administrator:} The
    super-user responsible for cross-tenant operations. Provisions
    new organisations and their email domains, tunes the platform's
    runtime configuration both deployment-wide and per organisation,
    invites the initial Educational Manager for each newly
    provisioned tenant, monitors system health, processing backlog
    and AI spend, and enforces platform-wide security and compliance
    policies. The IT Administrator holds the full platform permission
    set, since recovering a tenant from an operational failure
    requires it; in normal operation, however, course content and
    assessment configuration remain the responsibility of the
    Educational Manager and the Teachers of each tenant, and every
    administrative action taken inside a tenant is recorded in an
    append-only audit trail during its configured retention period.
\end{itemize}